\documentclass[11pt,titlepage,openright]{book}
\usepackage[utf8]{inputenc}
\usepackage[T1]{fontenc}
\usepackage[british]{babel}
\usepackage{graphicx}
\usepackage[dvipsnames]{xcolor}
\usepackage[sb]{libertine}
\usepackage[scale=0.9]{inconsolata}

\usepackage[sf]{titlesec}
\usepackage[round,sort,authoryear]{natbib}
%\usepackage[export]{adjustbox}

%% Math and algorithm support (added to align with the
%% Sjölund & Mair "Advice for Master's thesis reports" guidelines).
\usepackage{amsmath,amssymb,amsthm,mathtools}
\usepackage{algorithm}
\usepackage{algpseudocode}

\usepackage{
  marginnote, % Improved margin notes
  environ,
  ragged2e,   % Justified text in the margin notes
  url,        % For typesetting URLs
  listings,   % Code formatting
  hyperref,   % Links in PDF from TOC, refs, etc.
  lipsum,
  booktabs,
  changepage,
  float,
  dingbat
}
\renewcommand{\bfdefault}{bx}
\DeclareTextFontCommand\textsb{\libertineSB}

\usepackage[capitalise]{cleveref}
\Crefname{section}{\S}{\S\S}

\usepackage{enumitem}   % Improved lists 
% \setlist{noitemsep} % To make all lists compact
\setlist[itemize]{itemsep=0pt}
\setlist[itemize,1]{label=--}
\setlist[itemize,2]{label=\ensuremath{\triangleright}}
\setlist[itemize,3]{label=\RED{AVOID}}
% Description lists to use semibold labels
\setlist[description]{font=\libertineSB}

\usepackage[twoside,labelfont=sf]{caption}

\captionsetup{justification=raggedright,singlelinecheck=false}

\newcommand\myhrulefill[1]{\leavevmode\leaders\hrule height#1\hfill\kern0pt}
\DeclareCaptionFormat{FigFormat}{{\color{black}\myhrulefill{0.5pt}}\\#1#2#3}
\captionsetup[figure]{format=FigFormat}
\captionsetup[table]{format=FigFormat}

\DeclareCaptionFormat{LstFormat}{\textsf{Listing}~\arabic{chapter}.\arabic{listing}:#2#3}
\floatstyle{ruled}
\newfloat{listing}{thp}{lol}[chapter]
\floatname{listing}{Listing}
\captionsetup[listing]{format=LstFormat}

\NewEnviron{MarginNote}[1][0mm]{\marginnote{\footnotesize\justifying\BODY}[#1]}
\newcommand{\Footnote}[2][0mm]{\footnotemark\marginnote{\footnotesize$^{\arabic{footnote})}$~#2}[#1]}

\renewenvironment{figure*}[1][]{%
  \begin{figure}[#1]%
    % This feature is deprecated for now
    % \checkoddpage%
    % \ifoddpage%
    %   \begin{adjustwidth}{0cm}{-45mm}%%
    % \else%
    %   \begin{adjustwidth}{-45mm}{0cm}%%
    % \fi%
    }{%
    % \end{adjustwidth}%
  \end{figure}}

\renewenvironment{table*}[1][]{%
  \begin{table}[#1]%
    % This feature is deprecated for now
    % \checkoddpage%
    % \ifoddpage%
    %   \begin{adjustwidth}{0cm}{-45mm}%%
    % \else%
    %   \begin{adjustwidth}{-45mm}{0cm}%%
    % \fi%
    }{%
    % \end{adjustwidth}%
  \end{table}}

\renewenvironment{listing*}[1][]{%
  \begin{listing}[#1]%
    % This feature is deprecated for now
    % \checkoddpage%
    % \ifoddpage%
    %   \begin{adjustwidth}{0cm}{-45mm}%%
    % \else%
    %   \begin{adjustwidth}{-45mm}{0cm}%%
    % \fi%
    }{%
    % \end{adjustwidth}%
  \end{listing}}

%% == Code =======================================================
\lstnewenvironment{Code}[1][style=std]{\lstset{#1}}{}
\lstnewenvironment{Code_Numbered}[1][style=std,numbers=left]{\lstset{#1}}{}

\renewcommand{\c}[1]{\lstinline[style=std]@#1@}

\lstdefinestyle{std}{
  language=python,
  basicstyle=\small\tt\color{black},
  keywordstyle=\small\tt\bfseries,
  numberstyle=\footnotesize\sf\color{black},
  commentstyle=\small\color{black}\it,
  aboveskip=1ex,
  belowskip=1ex,
  tabsize=2,
  columns=fullflexible,
  xleftmargin=1ex,
  resetmargins=true,
  showstringspaces=false,
  morecomment=[l]{//},
  morecomment=[l]{--},
  morecomment=[s]{/*}{*/},
  escapeinside=@@,
  morekeywords={Frobies},
  moredelim=[is][\textit]{___}{___},
  moredelim=[is][\textbf]{__*}{*__},
  numberbychapter=true
}

\usepackage[activate={true,nocompatibility},final,tracking=true,kerning=true,spacing=true,factor=1100,stretch=10,shrink=10]{microtype}
\usepackage[paper=a4paper,text={13cm,24cm},marginparsep=5mm,marginparwidth=45mm,inner=20mm,twoside]{geometry}

\newcommand{\RED}[1]{\textcolor{red}{#1}}
\newcommand{\ie}{\emph{i.e.,}}
\newcommand{\eg}{\emph{e.g.,}}
\newcommand{\etal}{\emph{et~al.}}

\usepackage{pifont}
\newcommand{\Yes}{\ding{51}}
\newcommand{\No}{\ding{55}}

\renewcommand{\bfdefault}{b}
\clearpage{\pagestyle{empty}\cleardoublepage}

%% You can remove this line if you compile with --synctex=1 (see Makefile)
\synctex=1
\pagestyle{plain}

\begin{document}
\frontmatter
\title{Random Number Generator from Aggregated Cryptocurrency Prices \\ with an Application to Secure Password Generation}
\author{Yujia Liu \\ Uppsala University}
\date{\today}

\maketitle

\vspace*{3cm}
\section*{Abstract}

%% TODO: 论文摘要,约 200--300 词。建议结构:
%%   1) 问题背景与重要性(1--2 句);
%%   2) 现有方法的不足与本研究动机;
%%   3) 你提出的方法概述;
%%   4) 主要实验结果;
%%   5) 结论与影响。
\RED{Write the abstract last. It should mirror, in compressed form,
the Introduction and the Conclusion: setting, research questions,
method, key results, and takeaway.}


\tableofcontents
\listoffigures
\listoftables

\mainmatter

%% =====================================================================
%% Chapter 1 — Introduction
%% Guidelines (Introduction): 1) Motivation, 2) Brief literature review,
%% 3) Research questions.
%% =====================================================================
\chapter{Introduction}
\label{cha:introduction}

\section{Motivation}
\label{sec:motivation}

%% TODO (指南 Introduction §1):
%% Start with a high-level introduction to the area and then motivate
%% why the topic of the thesis is important.

\section{Brief Literature Review}
\label{sec:brief-lit-review}

%% TODO (指南 Introduction §2):
%% Briefly summarize the main existing methods and their shortcomings.

\section{Research Questions}
\label{sec:research-questions}

%% TODO (指南 Introduction §3):
%% In light of the motivation and existing knowledge, what are the
%% outstanding gaps that you intend to fill?

%% =====================================================================
%% Chapter 2 — Background
%% Guidelines (Background): 1) Literature review, 2) Background information.
%% =====================================================================
\chapter{Background}
\label{cha:background}

\section{Literature Review}
\label{sec:literature-review}

%% TODO (指南 Background §1):
%% This should not read like a shopping list, but instead synthesize
%% and compare different approaches related to your research questions.

\section{Background Information}
\label{sec:background-information}

%% TODO (指南 Background §2):
%% Provide sufficient detail so that one of your peers would be able
%% to understand the thesis. In other words, your thesis should be
%% self-contained for someone in the same program.

%% =====================================================================
%% Chapter 3 — Methods
%% Guidelines (Methods): 1) Focus on your contributions, 2) Use
%% mathematical formalism. (No further subsections prescribed.)
%% =====================================================================
\chapter{Methods}
\label{cha:methods}

%% TODO (指南 Methods §1, §2):
%% Describe your method in detail, focusing on the parts that are new
%% or different from previous work. Explain in what way your method
%% addresses your research questions. Use mathematical formalism where
%% appropriate; algorithms should preferably be typeset using the
%% algorithm environment.

%% =====================================================================
%% Chapter 4 — Experiments
%% Guidelines (Experiments): 1) Experimental setup, 2) Data,
%% 3) Baselines, 4) Results.
%% =====================================================================
\chapter{Experiments}
\label{cha:experiments}

\section{Experimental Setup}
\label{sec:experimental-setup}

%% TODO (指南 Experiments §1):
%% Ideally, it should be possible to relate each experiment to one of
%% the research questions, e.g., by proving or disproving a hypothesis,
%% studying the impact of a particular component using an ablation
%% study, or investigating the effect of varying a parameter.

\section{Data}
\label{sec:data}

%% TODO (指南 Experiments §2):
%% Describe how the data was acquired and what it contains. Also,
%% describe any preprocessing steps you perform.

\section{Baselines}
\label{sec:baselines}

%% TODO (指南 Experiments §3):
%% Explain which methods you compare against and what their setup is
%% (hyperparameters, etc.).

\section{Results}
\label{sec:results}

%% TODO (指南 Experiments §4):
%% Comment briefly on the results as you present them, highlighting
%% especially noteworthy things. In-depth discussions should be saved
%% for the discussion section. Strive to keep the presentation of a
%% given experiment relatively self-contained and focused.

%% =====================================================================
%% Chapter 5 — Discussion
%% Guidelines (Discussion): 1) Link results to research questions,
%% 2) Limitations and future work.
%% =====================================================================
\chapter{Discussion}
\label{cha:discussion}

\section{Linking Results to Research Questions}
\label{sec:linking-results-rq}

%% TODO (指南 Discussion §1):
%% Explain what the experiments tell you about your research questions.
%% If any of the experiments "failed" or did not behave as expected,
%% try to reason about potential explanations and propose additional
%% experiments that could be performed to lead the investigation further.

\section{Limitations and Future Work}
\label{sec:limitations-future-work}

%% TODO (指南 Discussion §2):
%% Describe the limitations of your work, i.e., under what assumptions
%% do your conclusions hold? What potentially relevant aspects were
%% outside the scope of your study? Possibly, suggest avenues for
%% future work building upon your study.

%% =====================================================================
%% Chapter 6 — Conclusion
%% Guidelines (Conclusion): 1) Recap the motivation and setting,
%% 2) Summarize your contributions in light of the experimental results.
%% =====================================================================
\chapter{Conclusion}
\label{cha:conclusion}

\section{Recap of the Motivation and Setting}
\label{sec:conclusion-recap}

%% TODO (指南 Conclusion §1):
%% The conclusion section is often quite similar to the abstract, but
%% longer. It should give a brief recap of the problem setting and the
%% research questions.

\section{Summary of Contributions}
\label{sec:summary-contributions}

%% TODO (指南 Conclusion §2):
%% Describe how the method you proposed manages to answer your
%% research questions.

%% =====================================================================
%% Bibliography
%% Guidelines (References §1): check style, no URLs/DOIs unless necessary,
%% prefer the published version over arXiv preprints when available.
%% =====================================================================
\bibliographystyle{plainnat}
\bibliography{references}

\end{document}

%% =====================================================================
%% Template syntax cheatsheet (everything below \end{document} is
%% ignored by LaTeX). Copy a block into the body when you need it.
%% These are taken from the original Wrigstad template.
%% =====================================================================

%% --- Figure (top of page) -------------------------------------------
\begin{figure}[t]
  \begin{center}
    \includegraphics[width=0.9\linewidth]{image.pdf}
  \end{center}
  \caption{Caption goes BELOW the figure.}
  \label{fig:example}
\end{figure}

%% --- Wide figure (extends into the margin) --------------------------
\begin{figure*}[t]
  \begin{center}
    \includegraphics[width=0.9\linewidth]{image.pdf}
  \end{center}
  \caption[Short caption for the LoF.]{Long caption for the body.}
  \label{fig:example-wide}
\end{figure*}

%% --- Table (booktabs; caption ABOVE the table) ----------------------
\begin{table}[b]
  \caption{Caption goes ABOVE the table.}
  \label{tab:example}
  \begin{center}
    \begin{tabular}{lll}
      \toprule
      $A$ & $B$ & $C$ \\
      \midrule
      $A$ & \multicolumn{2}{c}{$B-C$} \\
      $A$ & \Yes{} & \No{} \\
      \bottomrule
    \end{tabular}
  \end{center}
\end{table}

%% --- Inline code & code blocks --------------------------------------
Inline: \c{int inline = ++x;}.

\begin{Code}
public class HelloWorld<T> extends Something<T> {
  public static void main(String[] args) {
    System.out.println("Hello, world!");
  }
}
\end{Code}

\begin{Code_Numbered}
public class HelloWorld<T> extends Something<T> {
  public static void main(String[] args) {
    System.out.println("Hello, world!"); @\label{code:hello}@
  }
}
\end{Code_Numbered}

%% --- Code listing as a float (with caption) -------------------------
\begin{listing*}[t]
\begin{Code_Numbered}
public class HelloWorld<T> extends Something<T> {
  public static void main(String[] args) {
    System.out.println("Hello, world!");
  }
}
\end{Code_Numbered}
\caption{Code listings can also live inside a \c{listing} float.}
\label{lst:example}
\end{listing*}

%% --- Algorithm (uses the algorithm + algpseudocode packages) --------
\begin{algorithm}
  \caption{Example algorithm}
  \label{alg:example}
  \begin{algorithmic}[1]
    \Require Input $x$
    \Ensure Output $y$
    \State $y \gets f(x)$
    \State \Return $y$
  \end{algorithmic}
\end{algorithm}

%% --- Equation (use \eqref to cite it; equation is part of a sentence) -
The relation is given by
\begin{equation}
  y = f(x),
  \label{eq:example}
\end{equation}
which we refer back to using \eqref{eq:example}.

%% --- Verbatim (block + inline) --------------------------------------
\begin{verbatim}
plain monospaced text, no LaTeX commands processed
\end{verbatim}

Inline verbatim: \verb+\cref{code:hello}+ (use any matching delimiter
that does not appear in the verbatim text, e.g. \verb!...! or \verb|...|).

%% --- Margin note & footnote (margin) --------------------------------
\begin{MarginNote}[3cm]
  Margin note text. Use \texttt{\\-} to insert hyphens for line breaks
  in the narrow margin column.
\end{MarginNote}

Footnote text\Footnote[-1cm]{This footnote is rendered in the margin.}.

%% --- Lists ----------------------------------------------------------
\begin{itemize}
\item Standard itemize item.
  \begin{itemize}
  \item Nested item.
  \end{itemize}
\end{itemize}

\begin{itemize}[label=\leftpointright]
\item Custom label via enumitem.
\end{itemize}

\begin{enumerate}
\item Numbered item.
\item Another numbered item. \label{lst:two}
\end{enumerate}

\begin{enumerate}[resume]
\item Resumed numbering continues from above.
\end{enumerate}

\begin{description}
\item[Term] Definition or explanation.
\end{description}

%% --- Cross-references (cleveref) ------------------------------------
%% \Cref{...}  -- capitalised, use at the START of a sentence:
%%               "Chapter~3 ..." -> \Cref{cha:methods}
%% \cref{...}  -- lower-case, use mid-sentence:
%%               "as shown in chapter~3" -> as shown in \cref{cha:methods}
%% \eqref{...} -- equations only, gives "(N)":   see \eqref{eq:example}
%% \ref{...}   -- plain reference, no auto-prefix (rarely needed when
%%               cleveref is loaded).
%%
%% Targets work the same way for chapters / sections / subsections /
%% figures / tables / listings / algorithms / list items, e.g.:
%%   \cref{lst:two}      (item label inside an enumerate)
%%   \cref{code:hello}   (line label inside a Code_Numbered block)

%% --- Citations (natbib, Harvard) ------------------------------------
%% Active  (read it out loud):  \citet{key}        -> "Author (Year) shows ..."
%% Passive (not read out loud): \citep{key}        -> "... (Author, Year)."
%% Multiple keys:               \citep{key1,key2}
%% Year only:                   \citeyearpar{key}  -> "(Year)"
%% Author only:                 \citeauthor{key}   -> "Author"

%% --- Custom commands defined in the preamble ------------------------
\RED{Red highlight, used to flag TODO-style placeholders.}
\textsb{Semibold text (Libertine SB).}
\ie{} that is, used like this.
\eg{} for example, used like this.
\etal{} for "et al." in author lists.
\Yes{} \No{} dingbat check / cross marks (used in tables).

%%% Local Variables: ***
%%% mode: latex ***
%%% TeX-master: "main.tex"  ***
%%% ispell-local-dictionary: "british"  ***
%%% End: ***